%%%PAGE 169 rot=0 legibility=ok summary=铅笔随记：伏安法选器材与内外接、等时圆、光电门测a公式、双灯并串联（被划去）、电表非常规使用与改装、分压限流、磁电式电流表要点；右侧有试笔涂鸦
%%%BLOCK id=169-01 ch=10 sec=伏安法测电阻·选器材与内外接 kind=method exp=1 alt=none cont=none y=0.06-0.19
可以算出准确值 % 写在“精确与否”上方

选器材时 先看 \underline{精确与否} 后\unclear{度}大小阻 内外接 % “精确与否”下有下划线；“大小”与“阻”之间有一短横

阻值不\unclear{是}约数 % 写在“精确与否”下划线下方；“是”字形近“足”
%%%END
%%%BLOCK id=169-02 ch=03 sec=斜面·等时圆 kind=method alt=none cont=none y=0.17-0.27
%FIG p169_a 0.09 0.168 0.255 0.262
\notefig[0.28]{figures/p169_a.png}
从 $P$ 到两端时间相同

下方速度大\ 远
%%%END
%%%BLOCK id=169-03 ch=03 sec=牛顿第二定律实验·光电门测加速度 kind=note exp=1 alt=01 cont=none y=0.34-0.44
\[
mg=M\frac{d^2}{2x}\left(\frac{1}{t_2^2}-\frac{1}{t_1^2}\right)
\]
只能测 $t$ 求 $a$
%%%END
%%%BLOCK id=169-04 ch=10 sec=闭合电路·电源与灯泡U-I图线(双灯) kind=method exp=1 alt=none cont=none y=0.47-0.66
% 原稿此块被划去（每行均被横线划过）
双灯

并联\quad 满足同 $U$\quad $U=E-(I_1+I_2)r$\quad 找满足式子的直线
% 上行 E、r 下方各有一短横

串联\quad ：满足 同 $I$\quad $U=E-\unclear{I_2}\,(r_1+r_2+r)$\quad 找满足式子的直线
% “同 I”被圈出；“：”或为“∵”；式中括号前一项字形难辨；括号下另有一道长下划线

找 $U$ 等\unclear{于} % 小字，写在并联行与串联行之间，同样被划去
%%%END
%%%BLOCK id=169-05 ch=10 sec=电表改装·非常规使用 kind=summary exp=1 alt=none cont=none y=0.69-0.74
电表 $\to$ 非常规使用 $\to$ 改表 $\to$ 串联分压，并联分流

$U=IR_A$\quad $I=\dfrac{U}{R_V}$
%%%END
%%%BLOCK id=169-06 ch=10 sec=滑动变阻器·分压与限流 kind=method exp=1 alt=none cont=none y=0.74-0.84
\textbf{分压}
\begin{itemize}
\item $R$\unclear{全滑} $<\dfrac{1}{2}R_x$ 分压
\item 测量范围大
\item $U$-$I$ 曲线，示数校对 % CHECK: “U”后的字形近似“2”，据语境读作“-I”
\end{itemize}
其余用限流\qquad $U=\dfrac{E}{r+R'}R$
%%%END
%%%BLOCK id=169-07 ch=11 sec=安培力·磁电式电流表 kind=knowledge alt=10 cont=none y=0.87-0.99
磁电式电流表（$B$ 不是匀强 而均匀分布）

\underline{\unclear{转}向 同 $F_{\text{安}}$ 方向}，$B$ 不变

铝优点：不会被磁化\quad 减小对磁场影响

$\rho$ 小，电磁阻尼效果好 % “磁”字为后补，小字写在“电”“阻”之间的上方

运输时将接线柱用导线相连防损坏 % 写在页面右下；“运输时将接线柱用导线”下有划线，“相连防损坏”被框起
%%%END
%%%BLOCK id=169-08 ch=99 sec=涂鸦·试笔波浪线 kind=misc alt=none cont=none y=0.66-0.97
% 页面右侧大片试笔涂画：多条波浪线，其间写有两处英文 math；另在 y≈0.68 处有几笔零散折线
涂鸦：多条波浪线，其间写有 math、math
%%%END
%%%PAGEEND 169
